
% sciences/science/intro.tex

\setcounter{section}{-1}
\section{Introduction}

Solutions are proposed in Sections \ref{sec:mechanics}-\ref{sec:quantum} for exercises selected from \cite{TM, QM}.
Common notation is reviewed in Section \ref{sec:notation}.

\subsection{Notation}\label{sec:notation}

Denote the totally ordered field of real numbers by
\begin{align*}
	\left\langle \R, \R \times \R \xrightarrow{+} \R, 0, \R \times \R \xrightarrow{\cdot} \R, 1, \leq \right\rangle
\end{align*}
and the field of complex numbers as follows.
\begin{align*}
	\left\langle \C, \C \times \C \xrightarrow{+} \C, 0, \C \times \C \xrightarrow{\cdot} \C, 1 \right\rangle
\end{align*}
The field $\C$ of complex numbers carries the field involution $\C \xrightarrow{\overline{\bullet}} \C$ of complex conjugation.
Let $\C \times \C \xrightarrow{\braket{\bullet \mid \star}} \C$ denote the standard complex inner product on $\C$ given by $\braket{z \mid w} = \overline{z} w$, which is linear in $w$ and conjugate-linear in $z$.
The subset $\R_{\geq 0} = [0, \infty)$ of non-negative real numbers forms a totally ordered commutative semiring with multiplicative identity $1$.
Denote the totally ordered semiring of natural numbers as follows.
\begin{align*}
	\left\langle \N, \N \times \N \xrightarrow{+} \N, 0, \N \times \N \xrightarrow{\cdot} \N, 1, \leq \right\rangle
\end{align*}
For any natural $n \in \N$, denote the standard $n$-dimensional complex Hilbert space by
\begin{align*}
	\left\langle \C^n, \C^n \times \C^n \xrightarrow{+} \C^n, \vzero, \C \times \C^n \xrightarrow{\cdot} \C^n, \C^n \times \C^n \xrightarrow{\braket{\bullet \mid \star}} \C, \C^n \xrightarrow{\|\bullet\|} \R_{\geq 0} \right\rangle
\end{align*}
and the $n$-dimensional real Euclidean Hilbert space as follows.
\begin{align*}
	\left\langle \R^n, \R^n \times \R^n \xrightarrow{+} \R^n, \vzero, \R \times \R^n \xrightarrow{\cdot} \R^n, \R^n \times \R^n \xrightarrow{\cdot} \R, \R^n \xrightarrow{\|\bullet\|} \R_{\geq 0} \right\rangle
\end{align*}
Let $\Z_+$ be the totally ordered semigroup of positive integers.

\subsubsection{Euclidean plane}
The Euclidean plane is the $2$-dimensional real Euclidean Hilbert space $\R^2$.
Denote its standard basis unit vectors as follows.
\begin{align}\label{eqn:e2}
	\vx =
	\begin{pmatrix}
		1 \\ 0
	\end{pmatrix}
	\qquad
	\mbf{y} =
	\begin{pmatrix}
		0 \\ 1
	\end{pmatrix}
\end{align}
They form an ordered, orthonormal basis of the Euclidean plane $\R^2$.

\subsubsection{Euclidean space}
Denote the standard basis unit vectors in $\R^3$ as follows.
\begin{align}\label{eqn:e}
	\vx =
	\begin{pmatrix}
		1 \\ 0 \\ 0
	\end{pmatrix}
	\qquad \mbf{y} =
	\begin{pmatrix}
		0 \\ 1 \\ 0
	\end{pmatrix}
	\qquad \mbf{z} =
	\begin{pmatrix}
		0 \\ 0 \\ 1
	\end{pmatrix}
\end{align}
They form an orthonormal basis $e = (\vx, \mbf{y}, \mbf{z})$ of $\R^3$.
For any vector $\vec{n} \in \R^3$, denote its inner products with the standard basis vectors as follows.
\begin{align*}
	n_x = \vx \cdot \vec{n}
	\qquad
	n_y = \mbf{y} \cdot \vec{n}
	\qquad
	n_z = \mbf{z} \cdot \vec{n}
\end{align*}
Any triple of real numbers $\begin{pmatrix} n_x \\ n_y \\ n_z\end{pmatrix}$ uniquely determines the coordinate vector $\vec{n} \in \R^3$ satisfying the above equations.
Let $S^2 = \left\{\, \hat{n} \in \R^3 \mid \|\hat{n}\| = 1\right\}$ denote the $2$-dimensional unit sphere about the origin in $\R^3$.

\subsubsection{Outline}

In what follows, each section herein corresponds to a book whose exercises are attempted and each subsection corresponds to a lecture from the corresponding book.

